% !TEX program = xelatex
\documentclass[10pt,a4paper]{article}

% ======================== 基础包 ========================
\usepackage[top=0.3cm,bottom=0.3cm,left=1.2cm,right=1.2cm]{geometry}
\usepackage{ctex}
\usepackage{titlesec}
\usepackage{enumitem}
\usepackage{xcolor}
\usepackage{hyperref}
\usepackage{fontawesome5}
\usepackage{tikz}

% ======================== 颜色定义 ========================
\definecolor{primary}{HTML}{2B4C7E}
\definecolor{darktext}{HTML}{2D2D2D}
\definecolor{graytext}{HTML}{666666}
\definecolor{rulecolor}{HTML}{B0B0B0}

% ======================== 超链接 ========================
\hypersetup{
    colorlinks=true,
    linkcolor=primary,
    urlcolor=primary,
    pdfauthor={徐其栋},
    pdftitle={徐其栋 - 嵌入式/物联网工程师}
}

% ======================== 字体设置 ========================
\setCJKmainfont{SimHei}[BoldFont=SimHei]
\setmainfont{Segoe UI}[BoldFont=Segoe UI Bold]

% ======================== 段落和间距 ========================
\setlength{\parindent}{0pt}
\setlength{\parskip}{0pt}
\pagestyle{empty}

% ======================== 标题样式 ========================
\titleformat{\section}
    {\large\bfseries\color{primary}}
    {}
    {0em}
    {}
    [\vspace{-0.5em}{\color{rulecolor}\rule{\textwidth}{0.6pt}}\vspace{-0.2em}]

\titlespacing*{\section}{0pt}{0.35em}{0.15em}

% ======================== 自定义命令 ========================
\newcommand{\projectblock}[3]{%
    {\bfseries\color{darktext}#1}\hfill{\color{graytext}#2}\\%
    {\small\color{graytext}#3}%
    \vspace{0.2em}%
}

% ======================== 文档开始 ========================
\begin{document}

% ======================== 头部 ========================
\begin{center}
    {\fontsize{24}{28}\selectfont\bfseries\color{primary} 徐其栋}

    \vspace{0.1em}

    {\color{graytext}
    \faPhone\hspace{0.3em}17268891140
    \hspace{1.2em}
    \faEnvelope\hspace{0.3em}\href{mailto:xqd922@outlook.com}{xqd922@outlook.com}
    \hspace{1.2em}
    \faGlobe\hspace{0.3em}\href{https://1001.pp.ua}{个人主页}}

    \vspace{0.3em}

    {\color{darktext} 求职意向：\textbf{\color{primary}嵌入式开发工程师 / 物联网工程师}}
\end{center}

\vspace{0.2em}
{\color{rulecolor}\hrule height 0.8pt}

% ======================== 教育背景 ========================
\section{\faGraduationCap\hspace{0.4em}教育背景}

{\bfseries 滇西科技师范学院}\hfill{\bfseries 物联网工程\hspace{0.3em}|\hspace{0.3em}本科}

{\color{graytext} 2022.09 -- 2026.06}

\vspace{0.2em}

{\color{darktext}\textbf{主修课程：}C 语言、Python、单片机原理、嵌入式技术、计算机网络、传感器原理、数字电子技术}

% ======================== 专业技能 ========================
\section{\faCogs\hspace{0.4em}专业技能}

{\bfseries\color{primary}嵌入式 / 物联网：}\hspace{0.3em}{\color{darktext}STM32、ESP32、外设驱动调试、I2C、UART、MQTT、EMQX、蓝牙、Wi-Fi、FreeRTOS 初步}

\vspace{2pt}

{\bfseries\color{primary}编程与数据库：}\hspace{0.3em}{\color{darktext}C 语言、Python、C\#、MySQL、SQL 参数化查询}

\vspace{2pt}

{\bfseries\color{primary}开发工具：}\hspace{0.3em}{\color{darktext}Keil MDK、Visual Studio、Git、Linux、嘉立创 EDA}

% ======================== 实习经历 ========================
\section{\faBriefcase\hspace{0.4em}实习经历}

{\bfseries 云南航天工程物探检测股份有限公司}\hfill{\color{graytext} 2025.08 -- 2025.12}

{\small\color{graytext}技术助理\hspace{0.3em}|\hspace{0.3em}昆明\hspace{0.3em}|\hspace{0.3em}隧道 / 桥梁 / 路面检测设备}

\vspace{2pt}

\begin{itemize}[leftmargin=1.5em, itemsep=2pt, parsep=0pt, topsep=2pt]
    \item 基于 C/C++ 参与隧道、桥梁、路面检测设备底层软件开发，完成节点开发与传感器数据处理。
    \item 负责检测数据展示与功能交互开发，适配工程检测设备使用场景。
    \item 调试设备通信、排查程序 Bug，参与软件联调、性能优化与版本迭代，使用 Git 进行代码版本管理。
\end{itemize}

\vspace{0.15em}

% ======================== 项目经历 ========================
\section{\faProjectDiagram\hspace{0.4em}项目经历}

% --- 机器狗 ---
\projectblock{四足物联网机器狗控制系统}{2025.03 -- 2025.05}{技术栈：STM32F407ZE\hspace{0.3em}|\hspace{0.3em}PCA9685\hspace{0.3em}|\hspace{0.3em}ESP32\hspace{0.3em}|\hspace{0.3em}MPU6050\hspace{0.3em}|\hspace{0.3em}WS2812\hspace{0.3em}|\hspace{0.3em}I2C\hspace{0.3em}|\hspace{0.3em}UART\hspace{0.3em}|\hspace{0.3em}运动学逆解}

\begin{itemize}[leftmargin=1.5em, itemsep=2pt, parsep=0pt, topsep=2pt]
    \item 以 STM32F407ZE 为主控，结合 PCA9685 的 16 路 12 位 PWM 驱动 8 路舵机（四腿 $\times$ 2 关节），实现前进、后退、左右转、站立、蹲下、鞠躬、挥手等 9 类动作。
    \item 基于余弦定理与三角函数完成腿部运动学逆解，由足尖目标坐标 $(x, y)$ 反推大小腿舵机角度，并通过步进插值算法实现舵机平滑过渡，避免机械冲击。
    \item 搭建多通信链路：I2C 统一驱动 PCA9685 与 MPU6050；USART2 对接蓝牙实现网页遥控，UART4 对接 ESP32 完成语音识别，落地蓝牙、语音两种交互；WS2812 反馈系统状态。
\end{itemize}

\vspace{0.15em}

% --- 物联网云平台与数据库整合项目 ---
\projectblock{物联网云平台数据采集与三层架构管理系统}{2025.06 -- 2025.07}{技术栈：ESP32\hspace{0.3em}|\hspace{0.3em}MQTT\hspace{0.3em}|\hspace{0.3em}EMQX\hspace{0.3em}|\hspace{0.3em}C\#\hspace{0.3em}|\hspace{0.3em}.NET WinForms\hspace{0.3em}|\hspace{0.3em}MySQL\hspace{0.3em}|\hspace{0.3em}三层架构}

\begin{itemize}[leftmargin=1.5em, itemsep=2pt, parsep=0pt, topsep=2pt]
    \item \textbf{设备接入层}：ESP32 基于 MQTT 将传感器数据上报至 EMQX Broker，掌握 Topic、QoS、遗嘱消息等核心机制。
    \item \textbf{上位机 + 数据层}：独立实现 C\# 三层架构桌面端，DAL 参数化查询防注入，BLL 主键查重，形成"采集 → Broker → 上位机 → MySQL"端到端数据流。
\end{itemize}

\vspace{0.15em}

% --- 智能交互台灯 ---
\projectblock{STM32 智能交互台灯控制系统}{2026.03 -- 2026.05}{技术栈：STM32F103C8T6\hspace{0.3em}|\hspace{0.3em}光敏ADC\hspace{0.3em}|\hspace{0.3em}HC-SR04\hspace{0.3em}|\hspace{0.3em}SNR8016语音\hspace{0.3em}|\hspace{0.3em}HC-05蓝牙\hspace{0.3em}|\hspace{0.3em}PWM\hspace{0.3em}|\hspace{0.3em}OLED}

\begin{itemize}[leftmargin=1.5em, itemsep=2pt, parsep=0pt, topsep=2pt]
    \item 主控 STM32F103C8T6，集成光敏 ADC、超声波测距、红外感应、离线语音、蓝牙、PWM 调光、OLED 等 10 个模块，独立完成硬件选型与焊接调试。
    \item Keil C 模块化开发，模式状态机协调自动 / 手动 / 远程 / 语音四种交互方式。
\end{itemize}

\vspace{0.1em}

% ======================== 个人评价 ========================
\section{\faUser\hspace{0.4em}个人评价}

2026 届物联网工程本科毕业生，具备 STM32、ESP32 实战经验，熟悉 I2C、UART、MQTT、蓝牙通信；掌握 C\# + MySQL 上位机开发与 Git 版本管理，可快速进入嵌入式开发与测试工作。

\end{document}
